\begin{textsample}{POS Dim 5 – global\_south\_2024\_06 – Score 18.00 – global\_south\_2024\_06/109018199.txt}  \label{ex:f5_pos_222}
MEA has not used terms like ' Occupied Palestinian territory ' \&amp ; ' blatant disregard for \textbf{international} \textbf{law} ' in its statements, while trying to maintain a balancing \textbf{act} on Gaza.
New Delhi : India has given its stamp of approval to a joint statement by BRICS foreign ministers that is critical of Israel ' s actions in Gaza and its"disregard of \textbf{international} \textbf{law}".
The Ministry of External Affairs ( MEA ) has so far not used terms such as"occupied Palestinian territory"and"Israel ' s blatant disregard of \textbf{international} \textbf{law}"in its statements about the war, as it has tried to maintain a delicate balancing \textbf{act} between Israel and Palestine since last October.
When the war first broke out, Narendra Modi -- the first ever Indian PM to visit Israel -- condemned the 7 October attack by Hamas and expressed solidarity with Israel.
It set India apart from reactions in the ' Global South ', and hewed closer to the position of the US and other G7 countries at the time.
Five days later, the MEA clarified that @ @ @ @ @ @ @ @ @ @ need for a two-state solution.
Over the course of the war, India has \textbf{abstained} from key \textbf{resolutions} in the United Nations ( UN ) calling for an immediate ceasefire in Gaza, but backed Palestine ' s bid for \textbf{membership} at the UN last month.
New Delhi has provided 70 tons of humanitarian aid to Palestine while also facing pressure from Israel to designate Hamas as a terror group.
BRICS is a group of developing countries that comprises Brazil, Russia, India, China, South Africa and others. ' Occupied Palestinian Territory '"The Ministers expressed grave concern at the deterioration of the situation in the Occupied Palestinian Territory,"states paragraph 34 of the BRICS foreign ministers ' statement issued in Novgorod, Russia on Monday.
Rajiv Bhatia, former Indian ambassador to South Africa, told ThePrint :"The BRICS foreign ministers have used unusually harsh language to censure Israeli actions.
A key factor responsible for this noticeable development is the entry of Arab nations ( Saudi Arabia, Iran and UAE @ @ @ @ @ @ @ @ @ @ ) into the grouping."Notably, this was the first high-level BRICS meet since the entry of the three Middle Eastern nations.
Foreign ministers from all three countries attended the meeting, including Iran ' s \textbf{Acting} Minister of Foreign Affairs Ali Bagheri, who replaced the late Hossein Amir-Abdollahian.
Last November, Iran ' s then foreign minister wrote to foreign ministers from BRICS countries, urging their intervention in the"dire"Gaza situation.
He had also written to Shanghai \textbf{Cooperation} Organisation ( SCO ) head Zhang Ming with the same demand.
Palestine envoy welcomes BRICS statement The foreign ministers ' joint statement also"condemned the Israeli military operation in Rafah".
It expressed concern over the"high density of Palestinian civilians in this location, and the humanitarian catastrophic results due to the suspension of the Rafah crossing from the Palestinian side".
The statement, though briefly, also called for the release of Israeli hostages.
Speaking to ThePrint, Adnan Abu Al Haija, the Palestinian ambassador to India, @ @ @ @ @ @ @ @ @ @ and all parties of BRICS.
We are still looking for more pressure on Israel to think of peace, stop their \textbf{crimes} and to accept the two-state solution".
There was also a mention of South Africa ' s \textbf{case} against Israel at the \textbf{International} \textbf{Court} of Justice ( ICJ ) in the BRICS statement."They ( foreign ministers ) acknowledged the \textbf{provisional} \textbf{measures} of the \textbf{International} \textbf{Court} of Justice in the \textbf{legal} \textbf{proceedings} instituted by South Africa against Israel.
The Ministers expressed serious concern at Israel ' s continued blatant disregard of \textbf{international} \textbf{law}, the UN Charter, UN \textbf{resolutions} and \textbf{Court} orders,"the statement read.
After South Africa levelled \textbf{allegations} of \textbf{genocide} against Israel at the UN ' s top \textbf{court} last December, over a dozen countries, including Spain, Colombia and Turkey, expressed their support for the \textbf{case}.
Israel, however, has repeatedly rejected the charge.
In May, days before an Israeli air strike triggered a massive fire in a camp in Rafah, the ICJ issued a dramatic \textbf{ruling} in @ @ @ @ @ @ @ @ @ @ an emergency \textbf{measure} that was part of a larger \textbf{case} brought by South Africa.
Israel rejected the \textbf{ruling}."The invocation of the ICJ was likely inevitable, where South Africa spearheads the \textbf{accusations} levelled against Israel,"explained Israeli analyst and journalist Lev Aran.
He added :"Both South Africa and India are entering into the dharma of coalition ; however, from Israel ' s perspective, the electoral defeat of the \textbf{ruling} party in South Africa is welcomed news, and in the \textbf{case} of India, Israel hopes that India will serve as a voice of moderation in every forum of the Global South."

% matched lemmas: abstain, accusation, act, allegation, case, cooperation, court, crime, genocide, international, law, legal, measure, membership, proceeding, provisional, resolution, ruling
\end{textsample}
